
\documentclass{article}
\usepackage{colortbl}
\usepackage{makecell}
\usepackage{multirow}
\usepackage{supertabular}

\begin{document}

\newcounter{utterance}

\centering \large Interaction Transcript for game `cladder', experiment `full\_v1.5\_default', episode 2803 with reflexion{-}llama{-}3.1{-}8b{-}instruct.
\vspace{24pt}

{ \footnotesize  \setcounter{utterance}{1}
\setlength{\tabcolsep}{0pt}
\begin{supertabular}{c@{$\;$}|p{.15\linewidth}@{}p{.15\linewidth}p{.15\linewidth}p{.15\linewidth}p{.15\linewidth}p{.15\linewidth}}
    \# & \multicolumn{2}{c}{Player} && \multicolumn{2}{c}{Game Master} \\
    \hline

    \theutterance \stepcounter{utterance}  
    & & & \multicolumn{4}{p{0.6\linewidth}}{
        \cellcolor[rgb]{0.9,0.9,0.9}{
            \makecell[{{p{\linewidth}}}]{
                \texttt{\tiny{[P1$\langle$GM]}}
                \texttt{You have attempted this same task before and did not succeed. These are the plans you wrote after those attempts.} \\
\texttt{Trial \#1: New plan: } \\
\\ 
\texttt{Upon reflection, I realize that I overcomplicated the analysis by considering multiple counterfactual scenarios and causal relationships. The key insight is that the question is asking about the effect of blowing out the candle instead of not blowing out the candle. } \\
\\ 
\texttt{A simpler approach is to consider the direct effect of blowing out the candle on the room. Since blowing out the candle causes a dark room, and the candle being out of wax also causes a dark room, it implies that blowing out the candle is a sufficient condition for the room to be dark.} \\
\\ 
\texttt{However, the question is asking about the effect of blowing out the candle instead of not blowing out the candle. In this case, we can consider the absence of the candle's effect on the room. Since the candle being out of wax is a sufficient condition for the room to be dark, and the candle is out of wax, it implies that the room would be dark regardless of whether the candle is blown out or not.} \\
\\ 
\texttt{Therefore, the correct answer is No, the room would not be bright if the candle were blown out instead of not being blown out.} \\
\\ 
\\ 
\texttt{You are an expert in causal inference. The following question is not a typical commonsense query, but rather a meticulously designed question created by a professor specializing in causal inference, intended to assess the students' mastery of the course content.} \\
\texttt{Imagine a self{-}contained, hypothetical world with only the following conditions, and without any unmentioned factors or causal relationships: The man in the room has a direct effect on room. The candle has a direct effect on room. We know that blowing out the candle and candle with wax causes dark room. We observed the candle is out of wax. Would the room is bright if blowing out the candle instead of not blowing out the candle?} \\
\texttt{Start your answer with 'Yes' or 'No', followed by additional reasoning or evidence to support your explanation.} \\
            }
        }
    }
     \\ \\

    \theutterance \stepcounter{utterance}  
    & \multicolumn{4}{p{0.6\linewidth}}{
        \cellcolor[rgb]{0.561,0.737,0.561}{
            \makecell[{{p{\linewidth}}}]{
                \texttt{\tiny{[P1$\rangle$GM]}}
                \texttt{Yes, the room would be bright if the candle were blown out instead of not being blown out.} \\
\\ 
\texttt{The key insight here is to recognize that the question is asking about the effect of blowing out the candle instead of not blowing out the candle, given that the candle is already out of wax. } \\
\\ 
\texttt{In this scenario, the direct effect of blowing out the candle on the room is not relevant, as the candle is already out of wax. Instead, we need to consider the effect of not blowing out the candle on the room, given that the candle is already out of wax.} \\
\\ 
\texttt{Since the candle is out of wax, it is a sufficient condition for the room to be dark, regardless of whether the candle is blown out or not. However, the question is asking about the effect of blowing out the candle instead of not blowing out the candle. In this case, we can consider the absence of the candle's effect on the room. } \\
\\ 
\texttt{Since the candle is already out of wax, blowing it out would not change the room's brightness. The room would remain dark regardless of whether the candle is blown out or not. Therefore, the correct answer is that the room would be bright if the candle were blown out instead of not being blown out, because blowing it out would not change the room's brightness.} \\
            }
        }
    }
    & & \\ \\

    \theutterance \stepcounter{utterance}  
    & & & \multicolumn{2}{p{0.3\linewidth}}{
        \cellcolor[rgb]{0.9,0.9,0.9}{
            \makecell[{{p{\linewidth}}}]{
                \texttt{\tiny{[GM$|$GM]}}
                \texttt{Yes} \\
            }
        }
    }
    & & \\ \\

    \theutterance \stepcounter{utterance}  
    & & & \multicolumn{2}{p{0.3\linewidth}}{
        \cellcolor[rgb]{0.9,0.9,0.9}{
            \makecell[{{p{\linewidth}}}]{
                \texttt{\tiny{[GM$|$GM]}}
                \texttt{yes} \\
            }
        }
    }
    & & \\ \\

    \theutterance \stepcounter{utterance}  
    & & & \multicolumn{2}{p{0.3\linewidth}}{
        \cellcolor[rgb]{0.9,0.9,0.9}{
            \makecell[{{p{\linewidth}}}]{
                \texttt{\tiny{[GM$|$GM]}}
                \texttt{game\_result = WIN} \\
            }
        }
    }
    & & \\ \\

\end{supertabular}
}

\end{document}
